\subsection{Semisimple Modules over Tensor Products of Algebras}

In this subsection, we consider K-algebras \(A_{1}\) and \(A_{2}\) ; we denote the algebra \(A_{1} \otimes A_{2}\) by A.

PROPOSITION 1. --- Let \(M_{1}\) be an \(A_{1}\) -module and \(M_{2}\) an \(A_{2}\) -module, neither reduced to 0. If the module \(M = M_{1} \otimes M_{2}\) over the ring \(A = A_{1} \otimes A_{2}\) is simple (resp. isotypical, semisimple), then the \(A_{1}\) -module \(M_{1}\) and the \(A_{2}\) -module \(M_{2}\) are simple (resp. isotypical, semisimple).

Suppose that M is a semisimple A-module. Let \(N_{1}\) be an \(A_{1}\) -submodule of \(M_{1}\) . Denote the canonical image of \(N_{1} \otimes M_{2}\) in the A-module M by N. By Corollary 2 of VIII, p.~56, there exists an A-linear projector p in M with image N. By assumption, we have \(M_{2} \neq 0\) ; we can therefore choose an element m of \(M_{2}\) and a linear form \(\varphi\) on the K-vector space \(M_{2}\) such that \(\varphi(m) = 1\) (II, §7, No.~5, p.~302, Corollary 2). Let u be the mapping from \(M_{1}\) to M defined by \(u(m_{1}) = m_{1} \otimes m\) , and let v be the K-linear mapping from M to \(M_{1}\) characterized by \(v(m_{1} \otimes m_{2}) = \varphi(m_{2})m_{1}\) . Set \(q = v \circ p \circ u\) . The mapping \(q : M_{1} \to M_{1}\) is \(A_{1}\) -linear, its image is contained in \(N_{1}\) , and we have \(q(n) = n\) for every \(n \in N_{1}\) . Consequently, q is a projector in \(M_{1}\) with image \(N_{1}\) . We have proved that \(M_{1}\) is a semisimple \(A_{1}\) -module (VIII, p.~56, Corollary 2).

Suppose that M is simple and that \(M_{1}\) is the direct sum of two \(A_{1}\) -submodules \(M_{1}^{\prime}\) and \(M_{1}^{\prime\prime}\) . Set \(M^{\prime}=M_{1}^{\prime}\otimes M_{2}\) and \(M^{\prime\prime}=M_{1}^{\prime\prime}\otimes M_{2}\) ; then M is the direct sum of the A-submodules \(M'\) and \(M''\) . Since M is simple, \(M'\) or \(M''\) must be reduced to 0; as we have \(M_{2} \neq 0\) by assumption, we have \(M_{1}' = 0\) or \(M_{1}'' = 0\) (II, §3, No.~7, p.~256, Corollary 2). This proves that \(M_{1}\) is a simple \(A_{1}\) -module.

Now suppose that M is an isotypical A-module. Let S and T be simple \(A_{1}\) -submodules of \(M_{1}\) . The A-modules \(S \otimes M_{2}\) and \(T \otimes M_{2}\) can then be identified with nonzero submodules of M. They are therefore isotypical of the same type as M. By Remark VIII, p.~61, there exists a nonzero A-linear mapping \(f : S \otimes M_{2} \to T \otimes M_{2}\) . The mapping f is, in particular, \(A_{1}\) -linear. Since the \(A_{1}\) -modules \(S \otimes M_{2}\) and \(T \otimes M_{2}\) are nonzero and isotypical of type S and T, respectively, S and T are isomorphic. This proves that \(M_{1}\) is an isotypical \(A_{1}\) -module.

PROPOSITION 2. --- Let S be a simple module over the ring \(A = A_{1} \otimes A_{2}\) , finite-dimensional over K. For \(i \in \{1, 2\}\) , there exists a simple \(A_{i}\) -module \(S_{i}\) such that the \(A_{i}\) -module S is isotypical of type \(S_{i}\) . The A-module S is isomorphic to a quotient of the A-module \(S_{1} \otimes S_{2}\) .

Since S is a nonzero \(A_{1}\) -module that is finite-dimensional over K, it has finite length over \(A_{1}\) and there exist a simple left \(A_{1}\) -module \(S_{1}\) and a nonzero \(A_{1}\) -linear mapping from \(S_{1}\) to S. We endow \(\mathrm{M}_{2} = \mathrm{Hom}_{\mathrm{A}_{1}}(\mathrm{S}_{1}, \mathrm{S})\) with a left \(A_{2}\) -module structure defined by the external law \((a_{2}, u) \mapsto (a_{2})_{\mathrm{S}} \circ u\) . We have \(M_{2} \neq 0\) by construction, and \(M_{2}\) is finite-dimensional over K. We can therefore find a simple left \(A_{2}\) -module \(S_{2}\) and a nonzero \(A_{2}\) -linear mapping \(\varphi : S_{2} \to M_{2}\) . We define a nonzero A-linear mapping \(\psi\) from \(S_{1} \otimes S_{2}\) to S by

\[
\psi (s _ {1} \otimes s _ {2}) = \varphi (s _ {2}) (s _ {1})\tag{1}
\]

for every \(s_{1} \in S_{1}\) and \(s_{2} \in S_{2}\) . Since S is a simple A-module and \(\psi\) is not zero, \(\psi\) is surjective and S is isomorphic to a quotient of \(S_{1} \otimes S_{2}\) . For \(i \in \{1, 2\}\) , the \(A_{i}\) -module \(S_{1} \otimes S_{2}\) is isotypical of type \(S_{i}\) , and therefore so is the \(A_{i}\) -module S (VIII, p.~61, Proposition 2).

For any K-algebra B, we denote by \(\mathcal{S}_{\mathrm{K}}(\mathrm{B})\) the set of classes of simple (VIII, p.~51) left B-modules that are finite-dimensional over K.

THEOREM 1. --- Suppose that the field K is algebraically closed.

\begin{enumerate}
\def\labelenumi{\alph{enumi})}
\item
  Let \(M_{1}\) be an \(A_{1}\) -module and \(M_{2}\) an \(A_{2}\) -module, both simple (resp. semisimple) and finite-dimensional over K. Then \(M_{1} \otimes M_{2}\) is a simple (resp. semisimple) module over the ring \(A_{1} \otimes A_{2}\) and is finite-dimensional over K.
\item
  The mapping from \(\mathcal{S}_{\mathrm{K}}(\mathrm{A}_1)\times \mathcal{S}_{\mathrm{K}}(\mathrm{A}_2)\) to \(\mathcal{S}_{\mathrm{K}}(\mathrm{A}_1\otimes \mathrm{A}_2)\) that sends \((\operatorname {cl}(\mathrm{S}_1),\operatorname {cl}(\mathrm{S}_2))\) to \(\operatorname {cl}(\mathrm{S}_1\otimes \mathrm{S}_2)\) , where \(\mathrm{S}_1\) (resp. \(\mathrm{S}_2\) ) is a simple \(\mathrm{A}_1\) -module (resp. \(\mathrm{A}_2\) -module) that is finite-dimensional over \(\mathrm{K}\) , is bijective.
\end{enumerate}

To prove part a), it suffices to consider the case when \(M_{1}\) and \(M_{2}\) are simple. Let \(M'\) be an A-submodule of \(M = M_{1} \otimes M_{2}\) ; it is an \(A_{1}\) -submodule of \(M_{1} \otimes M_{2}\) , stable under the set of endomorphisms of the form \(1_{M_{1}} \otimes u\) , where u runs through the set of homotheties of the \(A_{2}\) -module \(M_{2}\) . Since the field K is algebraically closed, Schur's lemma (VIII, p.~47, Theorem 1) implies that the commutant \(\operatorname{End}_{\mathrm{A}_{1}}(\mathrm{M}_{1})\) of \(M_{1}\) is equal to K. By Corollary 2 of VIII, p.~63, the A-submodule \(M'\) of \(M_{1} \otimes M_{2}\) is of the form \(M_{1} \otimes M_{2}'\) , where \(M_{2}'\) is an \(A_{2}\) -submodule of \(M_{2}\) . We have assumed that \(M_{2}\) is simple; we therefore have \(M_{2}' = 0\) or \(M_{2}' = M_{2}\) , that is, \(M' = 0\) or \(M' = M\) . Consequently, M is simple.

If S is a simple module over \(A_{1} \otimes A_{2}\) that is finite-dimensional over K, then it follows from Proposition 2 and part a) that S is isomorphic to a module of the form \(S_{1} \otimes S_{2}\) , where \(S_{1}\) (resp. \(S_{2}\) ) is a simple \(A_{1}\) -module (resp. \(A_{2}\) -module). Moreover, as an \(A_{i}\) -module, S is isotypical of type \(S_{i}\) , so the class of \(S_{i}\) only depends on that of S. This proves part b).

Remarks. --- 1) Assertion a) of Theorem 1 is no longer true when the field K is not assumed algebraically closed. We can give examples (VIII, p.~225, Exercise 4) where \(M_{i}\) is a simple \(A_{i}\) -module that is finite-dimensional over K, for \(i \in \{1, 2\}\) , and where the A-module \(M_{1} \otimes M_{2}\) is not semisimple or is semisimple but not simple.

\begin{enumerate}
\def\labelenumi{\arabic{enumi})}
\setcounter{enumi}{1}
\tightlist
\item
  There exists a homomorphism \(\varphi\) from \(\mathrm{R_K(A_1)\otimes_{\mathbf{Z}} R_K(A_2)}\) to \(\mathrm{R_K(A)}\) characterized by the relation \(\varphi ([\mathrm{M}_1]\otimes [\mathrm{M}_2]) = [\mathrm{M}_1\otimes \mathrm{M}_2]\) . This can be proved in the same way as Proposition 9 of VIII, p.~196. If the field \(K\) is algebraically closed, then \(\varphi\) is an isomorphism from \(\mathrm{R_K(A_1)\otimes_{\mathbf{Z}} R_K(A_2)}\) to \(\mathrm{R_K(A)}\) by Theorem 1, b) because for every K-algebra B, the \(\mathbf{Z}\) -module \(\mathrm{R_K(B)}\) is free with basis the family \(([\mathrm{S}])_{\mathrm{S}\in \mathcal{S}_{\mathrm{K}}(\mathrm{B})}\) (VIII, p.~195).
\end{enumerate}

